\begin{helper}
Let \(V\) be finite, and let \(\nu\) be a measure on
\(\mathcal P(V)\times\mathbb R^V\), equipped with the product of the discrete
activation sigma-algebra and the Borel priority sigma-algebra.
Suppose every activation atom \(F_A=\{(B,q):B=A\}\) has finite mass, and
for every \(A\subseteq V\) and \(t:V\to[0,1]\),
\[
\nu\{(B,q):B=A,\ 0\le q(x)\le t(x)\text{ for all }x\in V\}
=\nu(F_A)\operatorname{ofReal}\!\left(\prod_{x\in V}t(x)\right).
\]
If \(D\subseteq\mathbb R^V\) is Borel measurable and has Lebesgue volume
zero, then \(\nu\{(B,q):q\in D\}=0\).
\end{helper}
\begin{proof}
Write \(Q=[0,1]^V\) and \(C=\{(B,q):q\in Q\}\).
Fix an activation atom \(F_A\). The priority projection is measurable,
and \(D\cap Q\) is a measurable subset of \(Q\) with volume zero by
monotonicity. Applying
\noderef{SamplingFiniteProductLowerOrthantEqualityMeasurable}
with event \(F_A\), priority projection, mass \(\nu(F_A)\), and set
\(D\cap Q\), gives
\[
\nu(F_A\cap\{q\in D\}\cap C)=\nu(F_A)\operatorname{vol}(D\cap Q)=0.
\]
The extra intersection with \(C\) in that theorem does not change this set.

The rectangle law at the constant threshold \(1\) gives
\(\nu(F_A\cap C)=\nu(F_A)\). The set \(F_A\cap C\) is measurable and
has finite mass. Subtracting its mass from the mass of \(F_A\) therefore
gives \(\nu(F_A\setminus(F_A\cap C))=0\).
Every point of \(F_A\cap\{q\in D\}\) either belongs to \(C\), in which
case it lies in the first null set, or does not, in which case it lies in
the second null set. Subadditivity and monotonicity show that
\(\nu(F_A\cap\{q\in D\})=0\).
Finally, the inverse image \(\{(B,q):q\in D\}\) is the union of these
sets over the finitely many activation atoms \(A\subseteq V\).
Finite subadditivity proves its nullity.
\end{proof}
